\section{手工构造 RNN：一个玩具示例}

为了深入理解 RNN 的前向传播过程，讲者通过一个手工设计的玩具示例，展示了如何构造权重矩阵来实现特定功能。

\subsection{任务定义}

\textbf{目标}：给定一个 0 和 1 的序列，当连续出现两个 1 时输出 1，否则输出 0。

\begin{figure}[H]
\centering
\includegraphics[width=0.95\textwidth]{slides-images/slide-011.jpg}
\caption{玩具示例任务：检测连续的两个 1。输入序列 $[0,1,0,1,1,1,0,1]$，对应输出 $[0,0,0,0,1,1,0,0]$\protect\footnotemark}
\end{figure}
\footnotetext{Slides 第12页。}

\subsection{隐藏状态设计}

为了完成这个任务，隐藏状态需要记录两条信息：\textbf{当前输入值}和\textbf{上一步的输入值}。讲者将隐藏状态设计为三维向量 $h_t = [\text{current}, \text{previous}, 1]^T$，其中第三个维度固定为 1（用于输出阶段的偏置计算）。

\begin{figure}[H]
\centering
\includegraphics[width=0.95\textwidth]{slides-images/slide-012.jpg}
\caption{隐藏状态设计：三维向量分别记录当前值、前一步的值和常数 1\protect\footnotemark}
\end{figure}
\footnotetext{Slides 第13页。}

\subsection{权重矩阵构造}

\textbf{输入权重 $W_{xh}$}：

$$W_{xh} = \begin{bmatrix} 1 \\ 0 \\ 0 \end{bmatrix}$$

当 $x_t = 1$ 时，$W_{xh} x_t = [1, 0, 0]^T$（记录当前值为 1）；当 $x_t = 0$ 时，结果为零向量。

\textbf{隐藏状态权重 $W_{hh}$}：

$$W_{hh} = \begin{bmatrix} 0 & 0 & 0 \\ 1 & 0 & 0 \\ 0 & 0 & 1 \end{bmatrix}$$

\begin{itemize}
    \item 第一行全零：当前值完全由 $W_{xh} x_t$ 决定
    \item 第二行 $[1,0,0]$：将上一步的``当前值''复制为这一步的``前一步值''
    \item 第三行 $[0,0,1]$：保持常数 1 不变
\end{itemize}

\textbf{输出权重 $W_{hy}$}：

$$W_{hy} = \begin{bmatrix} 1 & 1 & -1 \end{bmatrix}$$

输出 $y_t = \text{ReLU}(W_{hy} h_t) = \text{ReLU}(\text{current} + \text{previous} - 1)$。只有当 current $= 1$ 且 previous $= 1$ 时，$y_t = 1$。

\begin{figure}[H]
\centering
\includegraphics[width=0.95\textwidth]{slides-images/slide-015.jpg}
\caption{手工构造的完整 RNN 前向传播过程：通过精心设计的权重矩阵，RNN 能够正确检测连续的两个 1\protect\footnotemark}
\end{figure}
\footnotetext{Slides 第16页。}

\begin{knowledgebox}{从手工构造到梯度学习}
这个玩具示例的目的是帮助理解 RNN 的计算过程。在实际应用中：
\begin{itemize}
    \item 我们\textbf{不需要}手工设计权重矩阵——通过梯度下降自动学习
    \item 隐藏状态的维度通常远大于 3（可以是数百到数千维）
    \item 隐藏状态中学到的``特征''往往是高度抽象且不易解释的
\end{itemize}
\end{knowledgebox}

\subsection{本章小结}

通过手工构造 RNN 的权重矩阵，我们直观地看到了 RNN 如何通过隐藏状态在时间步之间传递信息。两组权重矩阵分别负责``处理当前输入''和``继承历史信息''，它们的加和经过激活函数后形成新的隐藏状态。

%% ========== RNN 训练 ==========

